\documentclass[a4paper]{cas-sc}
\usepackage[numbers]{natbib}
\usepackage[ruled,vlined,linesnumbered]{algorithm2e}
\DontPrintSemicolon
\SetAlgoVlined
\SetKwInput{KwIn}{Input}
\SetKwInput{KwOut}{Output}
\SetAlCapFnt{\small\bfseries}
\SetAlCapNameFnt{\small\bfseries}
\usepackage{amsmath,amssymb,graphicx,booktabs}
% Journal trim, type scale and column grid measured from the reference issue.
\geometry{paperwidth=210mm,paperheight=280mm,left=37.59bp,right=37.59bp,top=54.5bp,bottom=44bp,headheight=9bp,headsep=11bp,footskip=20bp}
\setlength{\columnsep}{18bp}
\setlength{\parindent}{12bp}
\renewcommand{\bibfont}{\fontsize{6.4}{8.1}\selectfont}
\setlength{\bibsep}{1.8bp plus .4bp minus .25bp}
\setlength{\emergencystretch}{1em}
\renewcommand{\normalsize}{\fontsize{8}{10.15}\selectfont}
\renewcommand{\small}{\fontsize{7.1}{8.7}\selectfont}
\renewcommand{\footnotesize}{\fontsize{6.4}{8.0}\selectfont}
\makeatletter
\@ifundefined{T@TU}{\input{tuenc.def}}{}
\DeclareFontFamily{TU}{JournalCharis}{}
\DeclareFontShape{TU}{JournalCharis}{m}{n}{<->"[fonts/CharisSIL_400Regular.ttf]:mapping=tex-text"}{}
\DeclareFontShape{TU}{JournalCharis}{m}{it}{<->"[fonts/CharisSIL_400Regular_Italic.ttf]:mapping=tex-text"}{}
\DeclareFontShape{TU}{JournalCharis}{b}{n}{<->"[fonts/CharisSIL_700Bold.ttf]:mapping=tex-text"}{}
\DeclareFontShape{TU}{JournalCharis}{b}{it}{<->"[fonts/CharisSIL_700Bold_Italic.ttf]:mapping=tex-text"}{}
\DeclareFontShape{TU}{JournalCharis}{bx}{n}{<->ssub*JournalCharis/b/n}{}
\DeclareFontShape{TU}{JournalCharis}{bx}{it}{<->ssub*JournalCharis/b/it}{}
\renewcommand{\encodingdefault}{TU}\renewcommand{\rmdefault}{JournalCharis}
\renewcommand{\sfdefault}{JournalCharis}\renewcommand{\ttdefault}{JournalCharis}
\makeatother
\setlength{\textfloatsep}{8bp plus 2bp minus 2bp}
\setlength{\dbltextfloatsep}{8bp plus 2bp minus 2bp}
\setlength{\floatsep}{6bp plus 2bp minus 1bp}
\setlength{\dblfloatsep}{7bp plus 2bp minus 1bp}
\renewcommand{\topfraction}{.93}\renewcommand{\dbltopfraction}{.93}
\renewcommand{\textfraction}{.04}\renewcommand{\floatpagefraction}{.72}\renewcommand{\dblfloatpagefraction}{.72}
\setcounter{topnumber}{3}\setcounter{dbltopnumber}{2}
\makeatletter
\renewcommand{\section}{\@startsection{section}{1}{\z@}{-10bp plus -1bp minus -1bp}{4bp}{\normalfont\normalsize\bfseries}}
\renewcommand{\subsection}{\@startsection{subsection}{2}{\z@}{-8bp plus -1bp minus -1bp}{4bp}{\normalfont\normalsize\itshape}}
\renewcommand{\subsubsection}{\@startsection{subsubsection}{3}{\z@}{-6bp plus -1bp minus -1bp}{2bp}{\normalfont\normalsize\itshape}}
\renewcommand{\@seccntformat}[1]{\csname the#1\endcsname.\enspace}
\renewcommand{\figurename}{Fig.}
\newcommand{\journalcaption}[2]{\par\vskip3bp\noindent\parbox{\linewidth}{\normalfont\fontsize{6.8}{8.4}\selectfont\textbf{#1.}\enspace#2\par}\par}
\newcommand{\jtablecaption}[2]{\par\noindent\parbox{\linewidth}{\normalfont\fontsize{6.8}{8.4}\selectfont\textbf{#1}\par#2\par}\par\vskip2.7bp}
\renewenvironment{figure}[1][tbp]{\@float{figure}[#1]\let\@makecaption\journalcaption\centering}{\end@float}
\renewenvironment{figure*}[1][tp]{\@dblfloat{figure}[#1]\let\@makecaption\journalcaption\centering}{\end@dblfloat}
\renewenvironment{table}[1][tbp]{\@float{table}[#1]\let\@makecaption\jtablecaption\centering\small\setlength{\tabcolsep}{2.6bp}}{\end@float}
\renewenvironment{table*}[1][tp]{\@dblfloat{table}[#1]\let\@makecaption\jtablecaption\centering\small\setlength{\tabcolsep}{3bp}}{\end@dblfloat}
\newcommand{\TableHere}[1]{\def\@captype{table}\let\@makecaption\jtablecaption\caption{#1}}
\makeatother
\newcommand{\Eq}[2]{\par\addvspace{4bp}\refstepcounter{equation}\label{eq:#1}%
 \noindent\makebox[\linewidth][c]{$\displaystyle #2$}\makebox[0pt][r]{(\theequation)}\par\addvspace{4bp}}
\newlength{\problemindent}\setlength{\problemindent}{8bp}
\newlength{\problemop}\setlength{\problemop}{53bp}
\newcommand{\ProbRow}[3]{\par\addvspace{3bp}\refstepcounter{equation}\label{eq:#1}%
 \noindent\hspace*{\problemindent}\makebox[\problemop][l]{$\displaystyle #2$}%
 \makebox[0pt][l]{$\displaystyle #3$}\hfill(\theequation)\par}


\newcommand{\MeasuredAbstract}{In 144 scene--seed pairs the method delivered 142 independently accepted motions, including 36/36 forward--reverse articulated maneuvers. Its median cost was 0.43\% lower than global covariance adaptation over 139 jointly feasible pairs.}
\newcommand{\CarCPU}{0.44}
\newcommand{\DockCPU}{0.27}

\newcommand{\ArticleTitle}{RegionCovOpt: An Integrated Metaheuristic Optimizer \mbox{for Generic Robotic Motion Planning}}
\newcommand{\ArticleAbstract}{Robotic motion planners often construct a coarse route before refining it by numerical optimization. The route generator must then decide how many alternatives to supply and how precisely to resolve each one before handing it to the local optimizer. This paper presents RegionCovOpt, a metaheuristic that integrates global exploration and local search under a single evaluation budget. Periodically reserved batches visit parameter regions in cyclic order and include an independent uniform draw. Other batches follow the region holding the best motion so far. Within a region, ranked trials guide local proposals and covariance adaptation. All trials use the same violation and cost criteria, so a promising motion discovered during exploration can redirect subsequent evaluations. The optimizer uses function values and requires no externally supplied trajectory. We study 48 parking, aerial flight, manipulator and articulated docking tasks, with matched controls, component ablations and held-out task perturbations. Endpoint-conditioned constructions account for much of the early feasibility. Relative to a parameter-matched global CMA control, median costs decrease by 28.49\% for flight and 27.98\% for manipulation, with little change for the two vehicle families. Against adaptive regional CMA, the held-out median change is $-0.035$\% and its scene-clustered interval spans zero. The coverage result concerns motions represented by the decoder. Code and experimental records are available at \href{https://github.com/libai1943/eaai}{github.com/libai1943/eaai}.}
\newcommand{\KeywordsBlock}{Robotic motion planning\par Metaheuristic optimization\par Regional covariance adaptation\par Trajectory optimization\par Collision avoidance}
% Typeset page furniture shared by the full-width title and two-column text.
\newcommand{\JournalName}{Engineering Applications of Artificial Intelligence}
\newcommand{\JournalURL}{https://www.sciencedirect.com/journal/engineering-applications-of-artificial-intelligence}
\definecolor{JournalBlue}{RGB}{0,128,172}
\hypersetup{colorlinks=true,citecolor=JournalBlue,linkcolor=JournalBlue,urlcolor=JournalBlue,
  pdfauthor={Bai Li},pdftitle={\ArticleTitle}}
\makeatletter
\def\ps@journal{\def\@oddhead{\fontsize{6.4}{8}\selectfont\itshape B. Li\hfil\JournalName}
  \let\@evenhead\@oddhead\def\@oddfoot{\hfil\fontsize{6.4}{8}\selectfont\thepage\hfil}\let\@evenfoot\@oddfoot}
\def\ps@journalfirst{\def\@oddhead{\hfil\fontsize{7}{9}\selectfont\color{JournalBlue}\JournalName\enspace---\enspace Research manuscript\hfil}\let\@evenhead\@oddhead\def\@oddfoot{}\let\@evenfoot\@oddfoot}
\makeatother
\newcommand{\EAFront}{%
  \par\noindent\begin{minipage}[t]{\textwidth}
  \setlength{\unitlength}{1bp}%
  \noindent\begin{picture}(520.1,82)
    \put(0,81.7){\rule{449.3bp}{.45bp}}
    \put(0,9.5){\includegraphics[width=59.55bp,height=65.35bp]{masthead/elsevier.png}}
    \put(72.3,9.5){\color[gray]{.9}\rule{377bp}{65.35bp}}
    \put(72.3,67.4){\makebox(377,6){\fontsize{8}{9}\selectfont Contents lists available at \href{https://www.sciencedirect.com/}{ScienceDirect}}}
    \put(72.3,39.6){\makebox(377,10){\fontsize{13.5}{16}\selectfont\JournalName}}
    \put(72.3,11.7){\makebox(377,6){\fontsize{7.5}{9}\selectfont journal homepage: \href{\JournalURL}{www.elsevier.com/locate/engappai}}}
    \put(460.6,9.5){\includegraphics[width=59.5bp,height=71.6bp]{masthead/journal.png}}
    \put(0,0){\rule{520.1bp}{2.35bp}}
  \end{picture}\par
  \vskip26bp\raggedright{\fontsize{13.5}{16.4}\selectfont\ArticleTitle\par}
  \vskip4.3bp{\fontsize{10.3}{12.2}\selectfont Bai Li\par}
  \vskip2.3bp{\fontsize{6.45}{8.2}\selectfont\itshape School of Information and Electronic Engineering, East China Normal University, Shanghai~200241, China\par}
  \vskip9bp\hrule height .4pt\vskip7bp
  \noindent\begin{minipage}[t]{.256\linewidth}
      {\fontsize{8}{10}\selectfont A R T I C L E\enspace I N F O\par}
      \vskip5.0bp\hrule height .45pt\vskip2.5bp
      {\fontsize{6.4}{8.2}\selectfont\textit{Keywords:}\par\KeywordsBlock\par}
  \end{minipage}\hfill
  \begin{minipage}[t]{.685\linewidth}
      {\fontsize{8}{10}\selectfont A B S T R A C T\par}
      \vskip5.0bp\hrule height .45pt\vskip2.5bp
      \fontsize{8}{10.1}\selectfont\rightskip=0pt\leftskip=0pt\parfillskip=0pt plus 1fil\ArticleAbstract\par
  \end{minipage}
  \par\vskip7bp\hrule height .45pt
  \vskip8bp\end{minipage}\par}


\begin{document}
\pagestyle{journal}\normalsize
\twocolumn[{\EAFront}]
\thispagestyle{journalfirst}
\begingroup\renewcommand{\thefootnote}{}\footnotetext{\textit{E-mail address:} \href{mailto:libai@zju.edu.cn}{libai@zju.edu.cn} (B. Li).}\endgroup
\section{Introduction}
Motion planning is a fundamental component of an autonomous robotic system. It determines how a robot reaches a prescribed state while respecting its motion capabilities and avoiding the surrounding obstacles. Sampling, graph search, reactive simulation, learning and numerical optimization offer different ways to accomplish this task. Their relative strengths become particularly important in a known, static but cluttered environment, where a feasible motion may leave considerable room for reducing maneuver time or control effort. This study is concerned with optimization of complete robotic motions under a finite computational budget.

An optimization-based planner expresses motion quality as an objective and enforces the robot's physical requirements through constraints. Continuous decision variables allow it to refine a trajectory beyond the resolution of a finite primitive library. Gradient-based methods exploit local sensitivity to make these refinements efficiently. However, the resulting trajectory depends on the initial guess, and nonconvex collision constraints can cause slow convergence or an unsuccessful solution process. A common planning architecture therefore first constructs a coarse route by sampling or search and then refines it through numerical optimization. The first stage must provide enough guidance for the second, while leaving enough computation for subsequent improvement.

The handoff between route search and refinement leaves two coupled decisions to the planner designer. How many route alternatives merit refinement, and how much geometric detail should the route generator resolve before the local optimizer takes over? A coarse route may leave a difficult maneuver for the second stage. Spending more time resolving that route leaves fewer evaluations for investigating alternatives. A collision discovered during refinement can call the route choice itself into question. In addition, derivative-based refinement needs sensitivities of the trajectory model and collision tests, which can be costly to obtain for discrete maneuvers, simulated dynamics or geometric penetration tests.

Function-value search offers an opportunity to address these issues within one iterative process. A population can investigate several motion alternatives and refine the useful ones without differentiating the motion evaluator. Yet removing derivatives also removes a powerful source of local guidance. Directly perturbing an entire control history can repeatedly miss the terminal state, and averaging coordinates across different maneuver constructions can destroy useful structure. Thus a derivative-free planner needs both a representation in which complete motions can be found and a search rule that uses its evaluations effectively.

RegionCovOpt makes each evaluation available to both global exploration and local improvement. Global exploration samples across the represented parameter domain. Local search uses the trial rankings to adapt a distribution in each region. Both operations evaluate complete motions, so the same measure of progress can guide the next allocation. A newly discovered motion with less violation, or lower cost at equal violation, immediately updates the best record. Its region receives the next unreserved batch if the record still leads. Periodically reserved batches revisit alternatives throughout the run. This gives the planner a way to redirect computation after a discovery, even late in the search.

The generic part of the method is the allocation and update rule. A robot-specific decoder and evaluator supply a motion and its rank for each trial. The optimizer uses function values and matrix operations without a third-party solver. We separate the contribution of motion construction from that of regional search by giving all optimizers the same decoder, initialization and evaluation budget. Matched global and regional CMA controls test the value of keeping distinct search distributions and of adding special proposals within them. The coverage analysis characterizes the archived cost under persistent independent exploration.

Section~2 reviews related work. Sections~3 and 4 formulate the planning problem and present RegionCovOpt. Section~5 reports the experiments, and Section~6 concludes.

\section{Related Works}
This section reviews the methods relevant to constructing feasible motions and improving their quality. We first discuss numerical optimal control and trajectory refinement, then examine the role of sampling and learned motion proposals. Finally, we consider derivative-free population search and regional allocation. This progression identifies which difficulties arise from motion representation and which remain within the optimizer.

\subsection{Optimal control formulations and trajectory refinement}
Paden et al.~\cite{survey} distinguished motion planning from feedback control in self-driving vehicles. Donald et al.~\cite{kino} formulated kinodynamic planning in terms of reachability under motion constraints. In a shooting formulation, integrating the controls eliminates the state equations, but the influence of a control on a later obstacle constraint must propagate through the intervening states. Hargraves and Paris~\cite{direct} treated states and controls as decision variables in direct collocation. Betts~\cite{betts} discussed the numerical trade-offs between trajectory optimization formulations. Kelly~\cite{kelly} showed how to construct and check a direct collocation transcription. The interior-point algorithm of W\"achter and Biegler~\cite{ipopt} exploits the derivative structure of the resulting nonlinear program.

An elastic band~\cite{elastic} deforms a path in response to its surrounding obstacles. CHOMP~\cite{chomp} instead updates a discretized trajectory using a covariant gradient. TrajOpt~\cite{trajopt} uses sequential convex approximations and checks the swept motion between adjacent configurations. Dong et al.~\cite{inference} represented motion planning as inference on a factor graph. Mukadam et al.~\cite{gpmp} used a continuous-time Gaussian process to interpolate between support states. These methods illustrate the value of a physically meaningful initial trajectory and of an obstacle model that supports local improvement.

The choice of local subproblem also matters. GuSTO~\cite{gusto} constructs a sequence of convex optimal control problems. Szmuk et al.~\cite{scp} treated state-triggered constraints by successive convexification. The convex feasible set algorithm of Liu et al.~\cite{cfs} restricts an update to a convex approximation of the feasible region. Zhang et al.~\cite{obca} used dual variables to impose separation between convex bodies. Gilbert et al.~\cite{rigid} computed the distance between convex objects from their support points. Van den Bergen~\cite{gk} developed a robust implementation of that geometric test. Accurate polygon geometry is consequently available to numerical optimization. The attraction of a function-value oracle here is that the same search can also evaluate a discrete maneuver choice or a nonsmooth geometry routine.

\subsection{Vehicle maneuvers and time parameterization}
Reeds and Shepp~\cite{reeds} characterized shortest bounded-curvature car paths with reversals in free space. In a cluttered workspace a connection between two poses must also admit the swept vehicle body and a feasible timing law. Ljungqvist et al.~\cite{tractor} addressed the additional articulation constraints in planning and following a path for a two-trailer system. Pham and Pham~\cite{topp} optimized the speed along a prescribed path by reachability analysis. In the present comparison every optimizer receives the same freedom to vary the geometric path and the same rules for assigning its duration. A difference in the reported maneuver time can therefore be attributed to the motion found by the search.

\subsection{Sampling-based motion planning}
Probabilistic roadmaps~\cite{prm} connect samples into a graph that can serve multiple queries. RRT-Connect~\cite{connect} grows trees from both boundaries to obtain a connection for a given task. LaValle and Kuffner~\cite{kino2} extended randomized tree growth to kinodynamic planning by propagating admissible controls. KPIECE~\cite{kpiece} uses the coverage of projected cells to decide where propagation should continue. OMPL~\cite{ompl} provides a common interface for such sampling-based planners.

Karaman and Frazzoli~\cite{rrt} established asymptotic optimality for specific graph and tree constructions. FMT*~\cite{fmt} propagates cost through a batch of samples. Informed RRT*~\cite{informed} restricts sampling once an incumbent bounds the region that can improve the solution. BIT*~\cite{bit} orders the search through an implicit graph using heuristic information. Li et al.~\cite{sst} analyzed sparse sampling-based kinodynamic planning. Their results concern the connectivity and cost of expanding graphs or trees. Our convergence question concerns the best complete motion generated by a fixed representation.

\subsection{Learning-based motion generation}
MPNet~\cite{mpnet} uses a trained network to propose motions and can invoke a classical planner when a connection fails. Kicki et al.~\cite{eaai1} trained a network to generate local car maneuvers from experience. Feh\'er et al.~\cite{eaai3} coupled reinforcement-learning path planning with closed-loop motion control and field tests. Carvalho et al.~\cite{diffusion} learned a multimodal trajectory prior and guided diffusion sampling with a task cost. In all four cases prior experience affects the proposals available for a new scene. RegionCovOpt instead obtains its prior information from the decoder, the specified robot model and the current map; it requires no training trajectories.

\subsection{Metaheuristic search and constraint handling}
STOMP~\cite{stomp} uses stochastic trajectory perturbations to improve a motion without computing its cost gradient. Particle swarm optimization~\cite{pso} uses personal and population records to direct particle updates. Differential evolution~\cite{de} constructs trials from population differences. The cross-entropy method~\cite{cem} fits a sampling distribution to selected trials. Wu~\cite{eaai4} combined constraint treatments with competition- and cooperation-driven differential evolution for uncertain parking. The present study treats known static geometry and tests every update strategy on the same motion families.

Covariance adaptation~\cite{cma} estimates dependencies between decision variables from the ranking of offspring. Hansen et al.~\cite{cma2} developed computationally efficient covariance updates. Hansen and Auger~\cite{cma3} related those updates to the design of continuous stochastic search. Pugh et al.~\cite{quality} distinguished solution quality from behavioral diversity. Vassiliades et al.~\cite{cvt} partitioned a descriptor space into a bounded number of archive cells. CMA-ME~\cite{cmame} coupled an archive with adaptive search distributions. Multi-Emitter MAP-Elites~\cite{memap} allocates work among heterogeneous emitters using a bandit mechanism. CMA-MAE~\cite{cmamae} adapts archive thresholds to balance objective improvement and behavioral coverage. These methods already establish the value of archives, multiple distributions and adaptive allocation. In the present formulation, regions identify a motion parameterization, candidate ranks enforce physical feasibility before cost, and the deliverable is a single complete motion. RegionCovOpt inherits covariance adaptation; its additional value must therefore be established against region-based CMA controls sharing that representation. Its region label identifies a part of the parameter space; different labels need not imply different homotopy classes.

Deb~\cite{deb} gave feasible candidates priority in selection. Stochastic ranking~\cite{runar} varied the influence of feasibility and objective value in ranking. Mesh adaptive direct search~\cite{mesh} and the pattern search of Lewis and Torczon~\cite{comp} offer derivative-free alternatives for constrained local improvement. Solis and Wets~\cite{solis} supplied conditions for convergence of random search. Our use of a reserved independent draw follows that line of reasoning. The distinctive empirical question is whether concentrating the remaining evaluations near a regional incumbent improves the cost of the motion that can actually be delivered.

\section{Problem formulation}
Let $z(t)$ denote the state of the robot and $u(t)$ its control at time $t$. The sets $\mathcal Z$ and $\mathcal U$ specify the admissible states and controls. The dynamics $f$ determine the rate of change of the state. The start state $z_{\mathrm s}$ and the goal state $z_{\mathrm g}$ are prescribed, and the terminal time $T$ is free. The occupied body at state $z$ is $\mathcal B(z)$. For the $j$th of $M$ static obstacles $\mathcal O_j$, let $g_j(z)$ denote signed clearance between $\mathcal B(z)$ and $\mathcal O_j$, positive when separated and negative under penetration. Let $\rho\ge0$ denote the required clearance.

A motion comprises a control history $u(\cdot)$, a state history $z(\cdot)$ and a terminal time $T>0$. The dynamics link these decision variables on $[0,T]$. The task cost $J$ can include a running cost, a terminal cost or the duration itself. The free-terminal-time problem is
\ProbRow{physical}{\mathop{\mathrm{minimize}}\limits_{\makebox[0pt]{$\scriptstyle u(\cdot),\,z(\cdot),\,T>0$}}}{J(u,z,T)}
\ProbRow{flow}{\mathrm{s.t.}}{\dot z(t)=f(z(t),u(t)),\quad\text{a.e. on }[0,T],}
\ProbRow{boundary}{}{(z(0),z(T))=(z_{\mathrm s},z_{\mathrm g}),}
\ProbRow{bounds}{}{(z(t),u(t))\in\mathcal Z\times\mathcal U,\; t\in[0,T],}
\ProbRow{gap}{}{g_j(z(t))\ge\rho,\; j\in\{1,\dots,M\},\ t\in[0,T].}
Equation~(\ref{eq:flow}) determines $z$ from the applied control and permits control switches. Equations~(\ref{eq:boundary})--(\ref{eq:gap}) impose the boundary states, physical limits and obstacle clearance, respectively. When the minimum is attained, the solution comprises the control $u^\star(t)$, the state $z^\star(t)$ and the terminal time $T^\star$. They jointly specify what the robot should execute and where it should be throughout the maneuver.

For the numerical search, write $\tau=(u(\cdot),z(\cdot),T)$ for a complete motion and abbreviate its cost as $J(\tau)$. A decoder $D$ maps a dimensionless parameter $\xi\in[0,1]^n$ to such a motion. It may impose endpoint conditions by construction or leave residuals to be tested. The implemented search partitions the parameter domain into two disjoint measurable regions $\mathcal R_b$, indexed by $b\in\mathcal I$. A trial is
\Eq{decoder}{\tau=D(\xi),\qquad \xi\in\bigcup_{b\in\mathcal I}\mathcal R_b.}
Let $V(\tau)$ be a nonnegative measure of violation, with zero denoting admission under the specified numerical tests. The admissible parameters in region $b$ and the represented optimal cost are
\Eq{feas}{\mathcal F_b=\{\xi\in\mathcal R_b:V(D(\xi))=0\}.}
\Eq{inf}{J_D^\star=\inf_{b\in\mathcal I}\;\inf_{\xi\in\mathcal F_b}J(D(\xi)).}
A decoder can omit an admissible physical maneuver. For this reason $J_D^\star$ is an infimum over represented motions. The independent physical replay in Section~5 provides an additional check of each returned motion.

The optimizer compares two candidates by the lexicographic order of
\Eq{rank}{R(\xi)=\bigl(V(D(\xi)),\,J(D(\xi))\bigr).}
A smaller violation wins; if violations are equal, a smaller cost wins. This ordering defines the incumbent before the first admission as well as after it. At budget exhaustion a feasible incumbent supplies its control and state profiles and terminal time. If its violation is positive, the procedure reports failure.
\section{Methodology}
RegionCovOpt chooses one region per batch and evaluates one motion per trial (Fig.~\ref{fig:workflow}). A trial that improves the regional record immediately changes the archive. A later trial in the same batch can already use that point as its center. At the end of the batch, the selected trials guide the update of the regional distribution. The global incumbent determines where the next unreserved batch runs.

\begin{figure*}[t]\centering\includegraphics[width=.89\textwidth]{../figures/scheme.pdf}
\caption{Integrated search in RegionCovOpt. A reserved batch selects a region cyclically and includes one uniform refresh when its slot is reached. Other batches follow the global incumbent. Each trial updates the records on improvement. Batch ranks adapt the selected distribution when the parent count is met.}\label{fig:workflow}\end{figure*}

\subsection{Initialization and regional records}
The optimizer draws $N_0$ initial points by Latin hypercube sampling of the unit box. Its evaluation budget is $B$, and the number of regions is $K=|\mathcal I|$. The evaluator returns a decoded motion and its rank from Eq.~(\ref{eq:rank}). Each visited region stores its best record, including the least-violating candidate before the first feasible motion. Unvisited regions start with empty archives. The best record across all regions is the global incumbent.

Each region stores a mean $m_b$, covariance $C_b$, scale $\sigma_b$, and two evolution paths. An affine transformation maps the selector coordinate to the unit interval of the local search space. The other coordinates already use that scale. This gives relative perturbations the same meaning across regions of different widths. A region with $n_b$ initial points starts at the mean of its best $\min(q_0,n_b)$ points, where $q_0$ limits the initial parent count. An empty region starts at the box center. Each covariance starts at the identity, each path at zero, and each scale at $\sigma_0$.

\subsection{Allocation and trial generation}
Let $H$ denote the reservation period in batches. Every $H$th batch, starting with the first, visits the next region in cyclic order. Each intervening batch goes to the region of the current global incumbent. Before admission, the incumbent is the point with the least violation, with cost breaking ties. After admission it is the motion with the lowest feasible cost. Batch capacity and reservation period set the granularity and minimum revisit rate. The observed ranks determine which region receives the intervening batches, and the regional distribution determines the scale and direction of local steps.

Each batch has capacity $\lambda$. The mean slot evaluates the regional mean. The poll slot perturbs a uniformly chosen coordinate of the regional record by a centered Gaussian increment with standard deviation $\eta_0+\eta_1\sigma_b$ when that record exists. The refresh slot of a reserved batch draws a uniform point in local coordinates. The optimizer fills every remaining slot with a correlated Gaussian draw. Its center is the regional record for an archive-eligible slot in $\mathcal J_A$ with a record available, and the regional mean otherwise. Section~5 gives the slot order and parameter values.

For a covariance factorization $C_b=Q\Lambda Q^\mathsf{T}$ and center $c$, the Gaussian trial is
\Eq{trial}{y=c+\sigma_b Q\Lambda^{1/2}\varepsilon,\qquad \varepsilon\sim\mathcal N(0,I).}
The optimizer reflects each trial coordinate at the unit boundaries and maps it back to the selected region. Decoding and ranking the resulting motion counts as one evaluation. A better rank updates the corresponding records. Once an archive contains a feasible motion, its cost can only stay the same or decrease.

\subsection{Distribution update}
Order a completed batch by increasing lexicographic rank. Let $\mu$ be the number of selected parents, $y_{i:\ell}$ the $i$th ranked point among the $\ell$ evaluated trials and $w_i$ the normalized log weights. The optimizer also updates the distribution from a truncated final batch when $\ell\geq\mu$. In local coordinates, omitting the region subscript for readability,
\Eq{weights}{w_i=\frac{\log(\mu+1/2)-\log i}{\sum_{j=1}^{\mu}[\log(\mu+1/2)-\log j]},\qquad i=1,\ldots,\mu.}
\Eq{eff}{\mu_{\mathrm{eff}}=\left(\sum_{i=1}^{\mu}w_i^2\right)^{-1}.}
\Eq{mean}{m^+=\sum_{i=1}^{\mu} w_i y_{i:\ell}.}
Write $d_i=(y_{i:\ell}-m)/\sigma$ for a selected displacement and $d=(m^+-m)/\sigma$ for the mean displacement. The first path accumulates whitened mean displacements. The second accumulates mean displacements in regional coordinates.
\Eq{ps}{p_\sigma^+=(1-c_\sigma)p_\sigma+\sqrt{c_\sigma(2-c_\sigma)\mu_{\mathrm{eff}}}\,C^{-1/2}d.}
\Eq{pc}{p_c^+=(1-c_c)p_c+h\sqrt{c_c(2-c_c)\mu_{\mathrm{eff}}}\,d.}
The gate $h$ suppresses the addition to $p_c$ when the normalized step-size path exceeds the threshold. For region-local update count $t$, dimension $n$ and the reference Gaussian length $\chi_n$ defined in Appendix A,
\Eq{gate}{h=\mathbf{1}\!\left\{\frac{\|p_\sigma^+\|}{\sqrt{1-(1-c_\sigma)^{2t}}}<(1.4+2/(n+1))\chi_n\right\}.}
\Eq{factor}{\gamma=1-c_1-c_\mu+(1-h)c_1c_c(2-c_c).}
\Eq{cov}{C^+=\gamma C+c_1p_c^+(p_c^+)^\mathsf{T}+c_\mu\sum_{i=1}^{\mu} w_i d_i d_i^\mathsf{T}.}
\Eq{scale}{\sigma^+=\operatorname{clip}_{[\sigma_{\min},\,\sigma_{\max}]}\!\left(\sigma\exp\!\left[\frac{c_\sigma}{d_\sigma}\!\left(\frac{\|p_\sigma^+\|}{\chi_n}-1\right)\right]\right).}
Appendix A gives the learning rates from Hansen and Ostermeier~\cite{cma} and the eigenvalue floor for the factorization. A symmetric eigendecomposition requires $O(n^3)$ operations per batch. Sampling and covariance adaptation add $O(\lambda n^2)$ operations, in addition to the robot-dependent cost of the oracle.

Algorithm~\ref{alg:region} gives the slot and batch updates. Appendix B uses the independent refreshes to bound the probability of missing a positive-mass near-optimal set.

\begin{algorithm}[tbp]
\small
\caption{RegionCovOpt}\label{alg:region}
\KwIn{Scene, boundaries, maps $D$ and $R$, evaluation budget $B$, initial size $N_0$, regional proposal schedule, batch capacity $\lambda$, parent count $\mu$, reservation period $H$.}
\KwOut{Best admitted control and state histories and their terminal time, or failure.}
Evaluate the $N_0$ stratified points; initialize regional records, distributions and the global incumbent\;
\For{each batch $e$ while budget remains}{
  \eIf{batch $e$ is reserved}{select the next region in cyclic order\;}{select the region of the global incumbent\;}
  Factor the selected covariance; start an empty batch $\mathcal L$\;
  \For{each slot in the scheduled batch while budget remains}{
    \uIf{the slot is the mean slot}{set the trial to the regional mean\;}
    \uElseIf{it is the poll slot and a regional record exists}{perturb one coordinate of that record\;}
    \uElseIf{it is the refresh slot of a reserved batch}{draw uniformly in local coordinates\;}
    \Else{draw a correlated Gaussian about the record if the slot is archive-eligible and the record exists, or about the mean otherwise\;}
    Reflect and map to the selected region; decode the motion and evaluate its rank\;
    Append the trial to $\mathcal L$; update both records if improved\;
  }
  \If{$|\mathcal L|\geq\mu$}{
    Rank $\mathcal L$ and update the selected mean, evolution paths, covariance and scale using Eqs.~(\ref{eq:mean})--(\ref{eq:scale})\;
  }
}
\If{the global incumbent is admissible}{\KwRet{its control, state and terminal time}\;}
\KwRet{failure and the least violation found}\;
\end{algorithm}


\section{Experiments and discussion}
We evaluate RegionCovOpt on four robot models. We compare the frequency of feasible returns, the costs of accepted motions, and the time each method takes. Matched controls and ablations then separate the contributions of motion construction, regional allocation, and trial generation.

\subsection{Experimental setup}
The benchmark contains twelve scenarios for each of four robot models. We tested fourteen settings in each scenario using seeds 4201, 5209, and 6203, for a total of 2,016 runs. Within each scene--seed pair, all methods use the same geometry, bounds, decoder, first $N_0$ evaluations, and total budget. Acceptance counts cover all runs. Paired cost comparisons use motions accepted for both methods. For timing, we ran three serial repeats on one logical processor.

\begin{center}\begin{minipage}{\linewidth}\centering\small
\TableHere{RegionCovOpt settings shared across scenes. Slot numbering starts at one.}\label{tab:settings}
\begin{tabular}{@{}p{.46\linewidth}p{.48\linewidth}@{}}\toprule
Quantity & Setting\\\midrule
Initial evaluations $N_0$; initialization parent limit $q_0$ & 24; 4\\
Regions $K$; reservation period $H$ & 2; every 4 batches, starting with the first\\
Batch capacity $\lambda$; selected parents $\mu$ & 12; 6\\
Initial scale $\sigma_0$; scale interval & 0.23; $[0.006,0.9]$\\
Coordinate-poll coefficients $(\eta_0,\eta_1)$ & $(0.005,0.4)$\\
Mean, coordinate poll, reserved refresh & First, second and third slot, respectively\\
Archive-centered Gaussian eligibility $\mathcal J_A$ & First five slots, subject to Section~4.2 precedence\\
Total evaluations $B$ & 6,000 (car); 960 (each other robot)\\
\bottomrule\end{tabular}\end{minipage}\end{center}

\subsubsection{Comparison methods and accounting}
We implemented eight controls in C++ under a common interface. They comprise global-best PSO, DE/rand/1/bin, diagonal Gaussian cross-entropy search, full-covariance CMA-ES, multistart finite-difference BFGS, random coordinate polling, uniform sampling, and the earlier regional poll-and-difference method. Decoding and scoring one candidate counts as one evaluation. All budgets include the same 24 initial candidates for a given scene and seed. Planning time includes path-word enumeration, scalar closure calculations, and the planner's feasibility replay. We measure independent verification separately.

PSO uses inertia 0.72984 and two gains of 1.49618. DE uses scale 0.8 and crossover probability 0.9. Cross-entropy search fits the best six of 24 trials with a coordinate standard-deviation floor of 0.025. CMA-ES uses 24 offspring, 12 log-weighted parents, full covariance and step-size adaptation. BFGS takes projected steps on $J+1000V$ with bounded finite differences and backtracking. Each trial used for a difference estimate or a backtracking step counts toward the budget. This BFGS comparison measures the performance of a gradient approximation in the search coordinates shared by all eight controls.

\subsection{Robot models, motion families and validation}
\subsubsection{Autonomous parking} The twelve public \href{https://github.com/libai1943/CommonParking}{CommonParking} cases use their original endpoint headings, vehicle geometry and obstacle polygons. The planner initializes its search from the map and boundary poses. The rear-axle bicycle model has wheelbase 2.8 m. Its state is $(x,y,\theta,v,\phi)$, and its controls are longitudinal acceleration $a=\dot v$ and steering rate $\omega=\dot\phi$. The parking decoder searches nine coordinates. The fourth coordinate divides the direct-connection and intermediate-pose regions at 0.35. Speed, acceleration, steering angle and steering rate have magnitude limits of 2.5 m/s, 1 m/s$^2$, 0.7 rad and 0.5 rad/s, respectively. The rigid rectangular body extends from $-0.929$ to 3.760 m longitudinally and from $-0.971$ to 0.971 m transversely. Dubins~\cite{dubins} characterized shortest forward-only paths under a curvature bound. Reeds and Shepp~\cite{reeds} characterized the corresponding paths when reversal is allowed. Our parking decoder connects oriented poses by signed circular and straight strokes, either directly or through an intermediate pose. Steering transitions occur at rest. The objective is complete maneuver time, including those transitions. The planner checks the rectangular body against a 0.01-m clearance threshold. Recursive midpoint bounds and a sufficient swept-projection shortcut check intersample motion. The independent replay samples at no more than 0.014 m of arc length, reconstructs the rectangle vertices, and checks final position and heading.

\subsubsection{Aerial flight and manipulator motion} The aerial model has three position coordinates and velocities. Its twelve maps contain two, three or four walls with openings at different lateral offsets and elevations. The flight model checks a spherical body of radius 0.22 m against the wall boxes. Each arm case has three joint coordinates and rates and two links of lengths 0.90 and 0.80 m. The arm model places four spheres of radius 0.13 m along each link and measures their clearance from the workspace boxes. The state comprises position and velocity in workspace coordinates for flight and in joint coordinates for the arm. Both models use the three accelerations as controls and integrate each held acceleration exactly. At interval length $\Delta t$ the common position update is
\Eq{spatial}{q_{k+1}=q_k+\Delta t\,v_k+\tfrac12\Delta t^2 a_k.}
\Eq{spatialv}{v_{k+1}=v_k+\Delta t\,a_k.}
The flight and arm decoders both use thirteen coordinates, with a split at 0.45 in the last coordinate. One region uses a velocity profile with zero endpoint velocity and prescribed integrated area, plus coordinate-wise shape changes whose integrated area is removed. The other uses internal knots, with bounded triangular velocity profiles between successive knots. In flight the knots can lie in the specified openings. In joint space they can lift the links, turn the base beyond the requested terminal azimuth, lower the arm and return. Mellinger and Kumar~\cite{snap} obtained quadrotor trajectories by optimizing a polynomial representation. Richter et al.~\cite{uav} exploited a piecewise-polynomial representation for flight through indoor environments. Zhou et al.~\cite{fast} incorporated flight-trajectory generation into replanning in unknown environments. Here the acceleration model allows both search regions to share the same propagation and acceptance tests. For both models the objective adds mean squared controls weighted by 0.01 and squared differences between consecutive controls weighted by 0.002 to the duration. Node clearance is 0.03 m. The violation measure adds the maximum bound excess, the largest dynamics defect, the peak clearance shortfall and its root-mean-square value across nodes. For trials that pass the node tests, open-loop replay subdivides each control interval into twenty steps and checks nonpenetration, state limits and terminal error below $10^{-3}$.

\subsubsection{Articulated reverse docking} The twelve loading bays differ in orientation, access and available staging space. Each task starts at rest in an aisle and ends at rest with the tractor and trailer aligned in a bay. Bay axes range from $44^\circ$ to $88^\circ$ relative to the map axis, on alternating sides of the aisle. All bays have a back wall and both side walls. The other features include a narrow mouth, staggered posts, adjacent occupied berths and a short or long apron. The start can have a nonzero heading. Both forward and reverse motion must occur, and the goal includes both final headings. Each scenario fixes the bay, obstacles, start, and goal.

The coordinates $(x,y)$ locate the tractor's rear axle. The angles $\theta$, $\psi$, and $\phi$ denote the tractor heading, trailer heading, and steering angle, respectively. The tractor wheelbase $L_{\mathrm t}$ and the rear hitch offset $\ell_{\mathrm h}$ are each 1.5 m. The distance $L_{\mathrm r}$ from the hitch to the trailer axle is 3 m. The articulated state augments the tractor position, heading, signed speed and steering angle with the trailer heading. The controls are $a$ and $\omega$, as for the car. With $v$ the signed speed, the heading equations are
\Eq{dhead}{\dot\theta=\frac{v\tan\phi}{L_{\mathrm t}}.}
\Eq{trhead}{\dot\psi=\frac{v\sin(\theta-\psi)-\ell_{\mathrm h}\dot\theta\cos(\theta-\psi)}{L_{\mathrm r}}.}
The tractor rectangle spans $[-0.25,1.8]$ m along the tractor axle frame, with half-width 0.90 m. The trailer rectangle spans $[-1.2,2.25]$ m along the trailer axle frame, with half-width 0.95 m. Both are tested against every oriented obstacle. Speed, acceleration, steering, steering rate and articulation magnitudes are limited to 1.5 m/s, 0.25 m/s$^2$, 0.7 rad, 0.5 rad/s and 1.2 rad.

The regions divide the first coordinate at 0.5. Seven unit-box coordinates choose the side and magnitude of a staging turn, the two bend amplitudes and their skews, and the staging extension. In T04, T08 and T11 a bollard lies ahead of the initial tractor nose. An eighth coordinate sets an initial backward withdrawal of 4.5--11 m, after which the vehicle changes to forward gear for its staging turn. The two regions use opposite staging turns. A bend has normalized arc coordinate $r\in[0,1]$, physical length $S$ and prescribed articulation $\alpha=\theta-\psi$. The amplitude $A$ and skew $k$ determine the articulation profile. Differentiating it with respect to signed distance $s$ and using the hitch equation gives the tractor curvature $\kappa$,
\Eq{alpha}{\alpha(r)=A\sin^2(\pi r)\,[1+k\sin(2\pi r)].}
\Eq{kappa}{\kappa(s)=\frac{\mathrm d\alpha/\mathrm ds+\sin(\alpha)/3}{1+\tfrac12\cos(\alpha)}.}
Both ends of the bend have zero articulation and steering. The integral of the derivative term over the bend vanishes because it is an exact differential and $\alpha(0)=\alpha(1)$. Thus for a prescribed heading change $\delta$ the length is set by
\Eq{length}{S=\frac{\delta}{\displaystyle\int_0^1 \frac{\sin(\alpha(r))}{3(1+\tfrac12\cos(\alpha(r)))}\,\mathrm dr}.}
The decoder constructs a forward bend from the start and an outward bend from the parked pose. A straight terminal stroke translates the outward bend along the bay axis so that both bend ends lie on the same staging line. Resolving their separation onto that line gives the remaining straight distance. Negative required lengths are violations. The rig follows the forward bend and advances to a stop, starting with a backward withdrawal in the three boxed-start tasks. It then reverses along the extension and outward bend into the berth.

Each stroke follows a half-cosine displacement profile with zero speed at both ends. Speed, acceleration, and steering-rate limits determine its duration, with a 1\% allowance. The objective sums the stroke times. The planner checks both rectangular bodies at 97 configurations per stroke against a 0.06-m projection-gap target. The independent verifier reconstructs speed and steering, then integrates the coupled headings from the original start using 400 RK4 steps per stroke. It carries the state across each gear change. Acceptance requires compliance with the physical limits, travel in both directions, terminal error below $10^{-3}$, and at least 0.03 m of body separation at the replay samples.

\subsection{Feasibility, route geometry and paired cost}
\begin{table}[htbp]\centering\small\caption{Independently accepted runs. Each family contributes twelve scenes and three seeds.}\begin{tabular}{lrrrrr}\toprule
Method & Car & Aerial & Arm & Trailer & Total\\\midrule
RegionCovOpt & 34 & 36 & 36 & 36 & 142\\
Particle swarm & 35 & 36 & 34 & 36 & 141\\
Full-covariance CMA-ES & 32 & 36 & 36 & 36 & 140\\
Differential evolution & 36 & 36 & 36 & 36 & 144\\
Gaussian cross-entropy & 32 & 36 & 36 & 36 & 140\\
FD--BFGS & 34 & 36 & 36 & 36 & 142\\
Coordinate polling & 29 & 36 & 35 & 36 & 136\\
Uniform sampling & 36 & 36 & 36 & 36 & 144\\
Poll/difference precursor & 36 & 36 & 36 & 36 & 144\\
\bottomrule\end{tabular}\end{table}
\begin{table}[htbp]\centering\small\caption{Paired primary costs relative to each control. Negative percentages favor RegionCovOpt. A successful motion outranks a failure in the win count.}\begin{tabular}{lrrrrr}\toprule
Control & Joint pairs & Median (\%) & Wins & Losses & Ties\\\midrule
Particle swarm & 139 & -0.23 & 95 & 45 & 4\\
Full-covariance CMA-ES & 139 & -0.43 & 105 & 36 & 3\\
Differential evolution & 142 & -0.42 & 97 & 44 & 3\\
Gaussian cross-entropy & 140 & -3.26 & 122 & 20 & 2\\
FD--BFGS & 140 & -0.32 & 96 & 48 & 0\\
Coordinate polling & 136 & -4.74 & 116 & 26 & 2\\
Uniform sampling & 142 & -6.27 & 125 & 19 & 0\\
Poll/difference precursor & 142 & -0.06 & 79 & 65 & 0\\
\bottomrule\end{tabular}\end{table}
\begin{table}[htbp]\centering\small\caption{Median paired cost change by robot family (percent), computed for runs accepted by both methods.}\begin{tabular}{lrrrr}\toprule
Control & Car & Aerial & Arm & Trailer\\\midrule
Particle swarm & +0.06 & -0.09 & -27.88 & -0.88\\
Full-covariance CMA-ES & -0.01 & -8.81 & -24.64 & +0.12\\
Differential evolution & +2.57 & -8.37 & -27.71 & -0.33\\
FD--BFGS & -3.24 & -24.24 & -0.35 & -0.03\\
\bottomrule\end{tabular}\end{table}

RegionCovOpt returned 142 independently accepted motions in 144 runs. It solved all twelve parking layouts at seed 4201. Its two failures were compact parallel parking P02 at seed 6203 and the two-sided pillar approach P09 at seed 5209. Differential evolution (DE) and uniform sampling each returned 144 accepted motions using the same decoder. With feasibility already common, paired cost reveals how effectively each method used the remaining evaluations.

We calculate the cost change for each pair of accepted motions as
\Eq{paired}{\Delta J=100\left(\frac{J_{\mathrm{RegionCovOpt}}}{J_{\mathrm{control}}}-1\right)\%.}
RegionCovOpt reduced median cost by 0.43\% relative to the original global CMA-ES control across 139 jointly feasible pairs. The largest reduction was 24.64\% in the arm tasks (36 pairs), followed by 8.81\% in flight. Against PSO, the arm reduction was 27.88\% across 34 pairs. Parking costs were close to those of CMA-ES ($-0.01$\%) and PSO ($+0.06$\%). In articulated docking, the median cost was 0.12\% higher than with CMA-ES. DE solved every parking run. On parking pairs accepted by both methods, RegionCovOpt cost 2.57\% more at the median.

Figure~\ref{fig:representative} shows three accepted motions per family at seed 4201. P02 and P09 move forward, reverse and finish in forward gear. P05 starts in reverse and backs into its final pose after a forward maneuver. All three articulated vehicles withdraw, move forward to align and reverse into the bay. The repository includes the complete scene plates and profiles.

\begin{figure*}[p]\centering
\includegraphics[width=.97\textwidth,height=.85\textheight,keepaspectratio]{../figures/representative.pdf}
\caption{Representative accepted motions at seed 4201. Rows show parking, flight, manipulation and articulated docking. Blue curves show RegionCovOpt and orange dotted curves show global CMA-ES. In the two vehicle rows, solid and dashed blue segments indicate forward and reverse travel. Diamonds mark gear changes and gray dashed rectangles show intermediate footprints. The planar axes use equal physical scales.}\label{fig:representative}
\end{figure*}

\subsection{State and control profiles}
For each scene, the profiles show the lowest-cost RegionCovOpt run that passed independent verification. We selected it from the three prescribed seeds, using the smaller seed to break ties, and reconstructed its controls and states from the initial state to the terminal time $T$. The two vehicle objectives equal $T$. The flight and arm objectives also include control effort and variation. The companion atlas lists the seed, $T$, and $J$ for all 48 scenes.

\begin{figure*}[t]\centering\includegraphics[width=.84\textwidth]{../figures/profile_ground.pdf}
\caption{Best accepted state and control histories for car parking P02 and articulated docking T04. Each column has its own time axis in seconds. Dotted vertical lines mark the onset of travel in the new gear. The angle $\theta$ is the tractor or car heading, $\psi$ the trailer heading and $\phi$ the steering angle.}\label{fig:groundprofiles}\end{figure*}
\begin{figure*}[t]\centering
\includegraphics[width=.84\textwidth]{../figures/profile_spatial.pdf}
\caption{Best accepted histories for aerial flight U04 and manipulator motion A04. The first two rows show configuration and velocity, and the third row shows the applied acceleration. Controls are held over each integration interval.}\label{fig:spatialprofiles}\end{figure*}

Figure~\ref{fig:groundprofiles} illustrates the effect of gear changes on elapsed time. P02 completes its maneuver in 20.99 s. Its signed speed returns to zero for steering changes at rest, so an interval with fixed position can still have nonzero steering rate. T04 starts by withdrawing, moves forward to prepare the approach and reverses into the bay. Forward travel begins at 9.65 s and the final reverse maneuver begins at 43.45 s. Its terminal time is 64.26 s. Figure~\ref{fig:spatialprofiles} shows the independently generated accelerations for U04 and A04. They reach their boundary states in 5.51 s and 2.13 s, respectively. Step changes in the acceleration are retained in the plots, and the state is propagated continuously across those changes.

\subsection{Component and representation ablations}
\begin{table}[htbp]\centering\small\caption{Component ablations. All rows compare the full method against a single change, except the last row which removes the entire special-proposal group.}\begin{tabular}{lrrrrr}\toprule
Variant & Accepted & Joint pairs & Median (\%) & Wins & Losses\\\midrule
Equal batch allocation & 144 & 142 & -0.16 & 97 & 47\\
Diagonal covariance & 144 & 142 & +0.00 & 64 & 77\\
Remove archive poll & 144 & 142 & +0.06 & 49 & 91\\
Remove reserved refresh & 141 & 140 & -0.00 & 74 & 65\\
Gaussian proposals only & 144 & 142 & -0.03 & 85 & 58\\
\bottomrule\end{tabular}\end{table}

Each variant receives the full budget, the same 24 initial members and the same motion representations as RegionCovOpt. Equal allocation spends four of every eight batches in each region. It reverses the alternation phase after four batches so that both regions retain a protected refresh. The diagonal variant removes learned off-diagonal covariances. Two other variants replace either the archive-coordinate poll or the independent reserved draw with the ordinary Gaussian proposal for that slot. The last ablation draws its entire batch about the distribution mean. All comparisons retain feasible alternatives found during the run.

Equal allocation returned two more feasible motions than RegionCovOpt. RegionCovOpt's median cost on their jointly accepted runs was 0.16\% lower. Removing the archive poll yielded 144 accepted motions. The full method cost 0.06\% more at the median on their jointly accepted runs. Dropping covariance cross-terms changed the median cost by less than 0.01\%. Sampling every trial from a Gaussian about the mean yielded 144 accepted motions, with RegionCovOpt 0.03\% lower on jointly accepted pairs. Removing the reserved refresh reduced the count to 141 and had little effect on median cost. The refresh supplies the independent draws in the coverage proof (Appendix B).

On P02 at seed 6203, RegionCovOpt spent 5,200 of 6,000 evaluations in the intermediate-pose region and ended with positive violation. Equal allocation spent 3,004 evaluations there and found its first feasible motion at evaluation 292. On P09 at seed 5209, equal allocation first reached feasibility at evaluation 5,066. Both cases illustrate the risk of concentrating work in the region with the lowest violation before any region has produced a feasible motion.

Direct control sampling yielded zero node-feasible motions in 96 runs. In this experiment, RegionCovOpt and DE each optimized six linearly interpolated control knots per channel and a free duration. We tested each method on all 48 scenes at seed 4201 with the original budgets. Forward integration determined the state history, and a penalty measured terminal error. This encoding increases the search dimension and leaves endpoint satisfaction to the search.

\subsection{Computation time}
\begin{table}[htbp]\centering\small\caption{Minimum process CPU time in seconds across three runs on one logical processor for scene 02 of each family.}\begin{tabular}{lrrrr}\toprule
Method & Car & Aerial & Arm & Trailer\\\midrule
RegionCovOpt & 0.44 & 0.19 & 0.23 & 0.27\\
Particle swarm & 0.38 & 0.19 & 0.20 & 0.28\\
Full-covariance CMA-ES & 0.38 & 0.19 & 0.22 & 0.28\\
Differential evolution & 0.42 & 0.16 & 0.20 & 0.27\\
Gaussian cross-entropy & 0.47 & 0.22 & 0.23 & 0.25\\
FD--BFGS & 0.66 & 0.20 & 0.25 & 0.25\\
Coordinate polling & 0.48 & 0.22 & 0.25 & 0.28\\
Uniform sampling & 0.77 & 0.12 & 0.09 & 0.27\\
\bottomrule\end{tabular}\end{table}

We timed case 02 of each robot on one logical processor and report the minimum process CPU time over three serial runs. Each repeat reproduced the original decision vector. RegionCovOpt took \CarCPU{} s for parking and \DockCPU{} s for articulated docking. PSO and CMA-ES each took 0.38 s for parking and 0.28 s for docking. On the arm case, the corresponding times were 0.23 s for RegionCovOpt, 0.20 s for PSO, and 0.22 s for CMA-ES.

\subsection{Matched controls and contribution attribution}

The matched controls use the same batch size (12), number of log-weighted parents (6), initial scale (0.23), scale limits, and learning rates. Global CMA fits one distribution. The two regional controls fit one distribution per normalized region and allocate batches either equally or by the incumbent/reservation rule. Each control initializes the mean from up to four top-ranked points in its domain and samples about that mean.

\begin{center}\begin{minipage}{\linewidth}\centering\small

\TableHere{Matched controls on the original 144 paired tasks and 128 held-out pairs. $\Delta J$ compares RegionCovOpt with the indicated control, using jointly accepted runs.}

\begin{tabular}{@{}lrrrr@{}}\toprule

Control & \multicolumn{2}{c}{Original} & \multicolumn{2}{c}{Held out}\\

 & Accepted & $\Delta J$ (\%) & Accepted & $\Delta J$ (\%)\\\midrule

RegionCovOpt & 142 & +0.000 & 128 & +0.000\\

Original CMA & 140 & -0.432 & 128 & -0.623\\

Matched global & 141 & -0.171 & 127 & -0.509\\

Regional, equal & 142 & -0.352 & 128 & -0.544\\

Regional, adaptive & 144 & -0.032 & 128 & -0.035\\

DE & 144 & -0.424 & 128 & -1.364\\

Uniform & 144 & -6.270 & 128 & -8.744\\

\bottomrule\end{tabular}\end{minipage}\end{center}

\begin{figure}[tb]
\centering
\includegraphics[width=\linewidth]{../figures/attribution.pdf}
\caption{Median paired cost change relative to matched global CMA on the original tasks (a) and adaptively allocated regional CMA on held-out tasks (b). Negative values favor RegionCovOpt. The panels use different horizontal scales. The bar in (b) gives the 95\% scene-cluster bootstrap interval. Only jointly accepted runs enter the comparison.}\label{fig:attribution}
\end{figure}


Figure~\ref{fig:attribution} shows where regional search reduced cost. Flight and manipulation had the largest median gains over matched global CMA. The original-panel median difference narrowed to $-0.032$\% against adaptive regional CMA. The regional controls incorporate both coordinate normalization and region-specific initial means, so this comparison measures their combined effect.

\subsection{Independent selection and held-out task perturbations}

We generated 12 validation tasks and 16 held-out tasks from separate sets of template indices. The perturbations varied the location and size of flight openings, arm goals and obstacles, parking starts and obstacles, and docking starts and aisle geometry. The repository contains the generator, seeds, and inputs.

The original rule led validation with 59 accepted motions in 60 runs. We evaluated five seeds per task and ranked variants first by acceptance count, then by median cost relative to the best accepted candidate for each pair. Removing polling, using diagonal covariance, balancing allocation before feasibility, and combining all three changes yielded 57, 57, 56, and 57 accepted motions, respectively. We fixed the selected rule before testing eight new seeds per held-out task. It and both regional controls returned 128/128 accepted motions. Matched global CMA returned 127/128.

Acceptance therefore favored the full rule on the validation tasks. The held-out cost comparisons below quantify the additional gain over regional CMA.

We estimated a separate 95\% percentile interval for each paired median using 4,000 bootstrap resamples. We resampled whole scenes within each robot family so that the seeds for a scene stayed together. These are marginal intervals for jointly accepted costs.

The held-out median cost change relative to matched global CMA was $-0.509$\% (95\% interval $[-18.400,-0.070]$\%).

The held-out median cost change relative to adaptive regional CMA was $-0.035$\% (95\% interval $[-0.103,0.011]$\%).

The shared initialization produced a motion that passed the planning checks in 125 of 144 original scene--seed pairs within 24 evaluations. This explains much of the early feasibility across optimizers. The budget curves track acceptance by the planner during search. The final cost comparisons use independently verified motions. The repository reports the evaluation count to first feasibility and the complete cost distributions.

\subsection{Representation scope and coverage}

The optimizer transfers across robots through the decoder and evaluator. The parking decoder builds bounded-curvature connections directly or through an intermediate pose. The flight decoder uses ordered rectangular openings, and the arm decoder uses velocity profiles or joint-space knots. The docking decoder uses prescribed gear sequences and aligns the tractor and trailer at the goal. New route branches, joint constraints, or gear sequences may require additional motion families.

At the 960-evaluation budget, each region receives ten independent refreshes after initialization. The probability of reaching a near-optimal motion in these draws depends on how much sampling probability the decoder assigns to that set. Appendix B expresses the bound in terms of that probability. The matched comparisons measure the costs reached at the experimental budget.


Both our C++ implementation of global CMA and pycma 4.5.0 reached objective error below $10^{-8}$ within 12,000 evaluations on all 12 rotated quadratic functions. The problems had dimensions 7, 9, and 13 and condition number $10^3$. Each implementation used its own boundary handling, initial means, and stopping rules. The supplementary records give the final error for each implementation.

\subsection{Planning and independently verified delivery time}

The extended timing test uses cases 02, 05 and 09 of each family, four search rules and three interleaved repeats per rule, giving 144 timed runs on one logical processor. Planning time includes process startup, decoding, scoring and all in-process replay. Independent acceptance is measured separately after reading the returned motion, with task data already loaded. Total wall time is their sum. The table gives median and 90th-percentile wall seconds for RegionCovOpt across the nine runs of each robot; the repository gives all controls and each case's best-of-three process CPU time.

\begin{center}\begin{minipage}{\linewidth}\centering\small

\TableHere{Planning and independent acceptance wall times in seconds. Each entry is median / 90th percentile.}

\begin{tabular}{@{}lrrr@{}}\toprule Robot & Planning & Acceptance & Total\\\midrule

Car & 0.517 / 0.628 & 0.004 / 0.007 & 0.530 / 0.632\\

Flight & 0.212 / 0.273 & 0.037 / 0.045 & 0.249 / 0.309\\

Arm & 0.271 / 0.333 & 0.017 / 0.018 & 0.287 / 0.352\\

Articulated & 0.406 / 0.426 & 0.406 / 0.413 & 0.809 / 0.830\\

\bottomrule\end{tabular}\end{minipage}\end{center}

Independent articulated replay is comparatively expensive because the Python verifier integrates every stroke with 400 RK4 steps and checks both bodies. Its cost should be distinguished from the compiled planner's runtime when assessing a delivered motion. The timings describe open-loop model-based planning and numerical acceptance; hardware execution is a separate validation question.


\section{Conclusion}
RegionCovOpt integrates exploration across parameter regions and local improvement within them. An exploratory trial can immediately improve the best record and redirect subsequent unreserved batches. The same rank and allocation rule apply to each robot through its decoder and evaluator.

Constructing motions that satisfy the endpoint conditions accounted for much of the early feasibility. Regional search then produced its largest cost gains in flight and manipulation. Adaptive regional CMA achieved similar costs to the full method on the held-out tasks. Extending the decoders to new map layouts would let us test whether these task-specific patterns persist.

\appendix
\section*{Appendix A. Regional update constants}
For search dimension $n$ and effective parent count $m=\mu_{\mathrm{eff}}$, the learning rates and damping factor are
\Eq{cc}{c_c=\frac{4+m/n}{n+4+2m/n}.}
\Eq{cs}{c_\sigma=\frac{m+2}{n+m+5}.}
\Eq{c1}{c_1=\frac{2}{(n+1.3)^2+m}.}
\Eq{cmu}{c_\mu=\min\left(1-c_1,\,\frac{2(m-2+1/m)}{(n+2)^2+m}\right).}
\Eq{damp}{d_\sigma=1+2\max\left(0,\sqrt{\frac{m-1}{n+1}}-1\right)+c_\sigma.}
\Eq{chi}{\chi_n=\sqrt n\left(1-\frac{1}{4n}+\frac{1}{21n^2}\right).}
Here $\chi_n$ approximates the expected norm of a standard Gaussian vector. Eigenvalues are floored at $10^{-14}$ when forming the square root and inverse square root of the covariance. The upper end of the selector coordinate in the lower region is clipped to $1-10^{-12}$ before affine mapping, to keep boundary points in their assigned region. If the last batch contains fewer than $\mu$ trials, its evaluations have already entered the archives and the distribution update is omitted.

\section*{Appendix B. Convergence of the best admissible cost}
Retain the parameter regions, numerical feasible sets and represented optimum of Section~3. Consider an infinite run with a deterministic decoder and rank. Let $\nu_b$ denote the law of a fresh local uniform draw mapped to region $b$, including the boundary handling of Appendix A. Assume that the reserved draws are mutually independent. Let $C_B$ denote the best admissible cost after $B$ evaluations, set to $+\infty$ before the first admission. For $\epsilon>0$, define the near-optimal subset and its probability by
\Eq{near}{A_{b,\epsilon}=\{\xi\in\mathcal F_b:J(D(\xi))<J_D^\star+\epsilon\}.}
\Eq{mass}{v_{b,\epsilon}=\nu_b(A_{b,\epsilon}).}

\textbf{Theorem 1.} Let $J_D^\star$ be finite. Suppose that for each $\epsilon>0$ at least one region $b\in\mathcal I$ has $v_{b,\epsilon}>0$. With one independent refresh in every reserved batch, the best admissible cost converges almost surely to $J_D^\star$. At budget $B=N_0+\lambda N$, after $N$ full batches, its upper-tail probability satisfies
\Eq{bound}{\Pr\{C_B\ge J_D^\star+\epsilon\}\le\prod_{b\in\mathcal I}(1-v_{b,\epsilon})^{\lfloor N/(HK)\rfloor}.}
A factor with zero exponent is interpreted as one.

\textit{Proof.} Each full block of $HK$ batches contains exactly $K$ reservations. Cyclic assignment sends one of these batches to each region. Consequently each region receives at least $\lfloor N/(HK)\rfloor$ independent refreshes among the first $N$ batches. A draw in $A_{b,\epsilon}$ has zero violation and cost below $J_D^\star+\epsilon$. The lexicographic acceptance rule retains that cost or a smaller one for the rest of the run. The event on the left of Eq.~(\ref{eq:bound}) therefore entails that every reserved draw missed its corresponding set. Independence gives the product bound. This step requires independence only for the reserved draws.

For each integer $j\ge1$, take $\epsilon=1/j$. By hypothesis the right-hand side of Eq.~(\ref{eq:bound}) tends to zero. The event that the archive never reaches this neighborhood is contained in every finite-budget miss event and hence has probability zero. Taking the countable intersection of the probability-one hit events shows that, almost surely, all these neighborhoods are eventually reached. The archive costs are nonincreasing and bounded below by $J_D^\star$ once finite. They therefore converge to $J_D^\star$. \hfill$\square$

The positive-mass assumption connects the proof to the engineering problem. For example, an admissible parameter with a feasible interior neighborhood and a continuous objective admits positive-probability draws arbitrarily close to its cost if the refresh law has positive density there. A terminal equality applied solely after forward integration can instead restrict admissible parameters to a set of zero volume. A maneuver omitted by the decoder cannot contribute to $J_D^\star$. The theorem concerns the represented admissible set and an ideal independent-draw process. A finite pseudorandom run and its numerical replay support separate empirical claims about a stated budget and acceptance tests.


\section*{Data and code availability}
The \href{https://github.com/libai1943/eaai}{public repository} provides the C++ source, all task inputs, the independent verifier and the run-level comparison and ablation records. The parking inputs derive from CommonParking revision \href{https://github.com/libai1943/CommonParking/tree/476ef8573d9e22578be5c3d42514a9a29f27d07a}{476ef857}. The 48-case state and control atlas identifies the selected run for each scene.
\begin{thebibliography}{99}
\fontsize{6.4}{8.1}\selectfont
\bibitem{survey} Paden, B., Cap, M., Yong, S.Z., Yershov, D., Frazzoli, E., 2016. A Survey of Motion Planning and Control Techniques for Self-Driving Urban Vehicles. IEEE Transactions on Intelligent Vehicles 1, 33-55. \href{https://doi.org/10.1109/tiv.2016.2578706}{doi:10.1109/tiv.2016.2578706}.
\bibitem{kino} Donald, B., Xavier, P., Canny, J., Reif, J., 1993. Kinodynamic motion planning. Journal of the ACM 40, 1048-1066. \href{https://doi.org/10.1145/174147.174150}{doi:10.1145/174147.174150}.
\bibitem{direct} Hargraves, C.R., Paris, S.W., 1987. Direct trajectory optimization using nonlinear programming and collocation. Journal of Guidance, Control, and Dynamics 10, 338-342. \href{https://doi.org/10.2514/3.20223}{doi:10.2514/3.20223}.
\bibitem{betts} Betts, J.T., 1998. Survey of Numerical Methods for Trajectory Optimization. Journal of Guidance, Control, and Dynamics 21, 193-207. \href{https://doi.org/10.2514/2.4231}{doi:10.2514/2.4231}.
\bibitem{kelly} Kelly, M., 2017. An Introduction to Trajectory Optimization: How to Do Your Own Direct Collocation. SIAM Review 59, 849-904. \href{https://doi.org/10.1137/16m1062569}{doi:10.1137/16m1062569}.
\bibitem{ipopt} Wächter, A., Biegler, L.T., 2005. On the implementation of an interior-point filter line-search algorithm for large-scale nonlinear programming. Mathematical Programming 106, 25-57. \href{https://doi.org/10.1007/s10107-004-0559-y}{doi:10.1007/s10107-004-0559-y}.
\bibitem{elastic} Quinlan, S., Khatib, O., 1993. Elastic bands: connecting path planning and control. [1993] Proceedings IEEE International Conference on Robotics and Automation, 802-807. \href{https://doi.org/10.1109/robot.1993.291936}{doi:10.1109/robot.1993.291936}.
\bibitem{chomp} Zucker, M., Ratliff, N., Dragan, A.D., Pivtoraiko, M., Klingensmith, M., Dellin, C.M., Bagnell, J.A., Srinivasa, S.S., 2013. CHOMP: Covariant Hamiltonian optimization for motion planning. The International Journal of Robotics Research 32, 1164-1193. \href{https://doi.org/10.1177/0278364913488805}{doi:10.1177/0278364913488805}.
\bibitem{trajopt} Schulman, J., Duan, Y., Ho, J., Lee, A., Awwal, I., Bradlow, H., Pan, J., Patil, S., Goldberg, K., Abbeel, P., 2014. Motion planning with sequential convex optimization and convex collision checking. The International Journal of Robotics Research 33, 1251-1270. \href{https://doi.org/10.1177/0278364914528132}{doi:10.1177/0278364914528132}.
\bibitem{inference} Dong, J., Mukadam, M., Dellaert, F., Boots, B., 2016. Motion Planning as Probabilistic Inference using Gaussian Processes and Factor Graphs. Robotics: Science and Systems XII. \href{https://doi.org/10.15607/rss.2016.xii.001}{doi:10.15607/rss.2016.xii.001}.
\bibitem{gpmp} Mukadam, M., Dong, J., Yan, X., Dellaert, F., Boots, B., 2018. Continuous-time Gaussian process motion planning via probabilistic inference. The International Journal of Robotics Research 37, 1319-1340. \href{https://doi.org/10.1177/0278364918790369}{doi:10.1177/0278364918790369}.
\bibitem{gusto} Bonalli, R., Cauligi, A., Bylard, A., Pavone, M., 2019. GuSTO: Guaranteed Sequential Trajectory optimization via Sequential Convex Programming. 2019 International Conference on Robotics and Automation (ICRA), 6741-6747. \href{https://doi.org/10.1109/icra.2019.8794205}{doi:10.1109/icra.2019.8794205}.
\bibitem{scp} Szmuk, M., Reynolds, T., Acikmese, B., Mesbahi, M., Carson, J.M., 2019. Successive Convexification for 6-DoF Powered Descent Guidance with Compound State-Triggered Constraints. AIAA Scitech 2019 Forum. \href{https://doi.org/10.2514/6.2019-0926}{doi:10.2514/6.2019-0926}.
\bibitem{cfs} Liu, C., Lin, C.Y., Tomizuka, M., 2018. The Convex Feasible Set Algorithm for Real Time Optimization in Motion Planning. SIAM Journal on Control and Optimization 56, 2712-2733. \href{https://doi.org/10.1137/16m1091460}{doi:10.1137/16m1091460}.
\bibitem{obca} Zhang, X., Liniger, A., Borrelli, F., 2021. Optimization-Based Collision Avoidance. IEEE Transactions on Control Systems Technology 29, 972-983. \href{https://doi.org/10.1109/tcst.2019.2949540}{doi:10.1109/tcst.2019.2949540}.
\bibitem{rigid} Gilbert, E.G., Johnson, D.W., Keerthi, S.S., 1988. A fast procedure for computing the distance between complex objects in three-dimensional space. IEEE Journal on Robotics and Automation 4, 193-203. \href{https://doi.org/10.1109/56.2083}{doi:10.1109/56.2083}.
\bibitem{gk} Bergen, G.V.D., 1999. A Fast and Robust GJK Implementation for Collision Detection of Convex Objects. Journal of Graphics Tools 4, 7-25. \href{https://doi.org/10.1080/10867651.1999.10487502}{doi:10.1080/10867651.1999.10487502}.
\bibitem{reeds} Reeds, J., Shepp, L., 1990. Optimal paths for a car that goes both forwards and backwards. Pacific Journal of Mathematics 145, 367-393. \href{https://doi.org/10.2140/pjm.1990.145.367}{doi:10.2140/pjm.1990.145.367}.
\bibitem{tractor} Ljungqvist, O., Evestedt, N., Axehill, D., Cirillo, M., Pettersson, H., 2019. A path planning and path-following control framework for a general 2-trailer with a car-like tractor. Journal of Field Robotics 36, 1345-1377. \href{https://doi.org/10.1002/rob.21908}{doi:10.1002/rob.21908}.
\bibitem{topp} Pham, H., Pham, Q.C., 2018. A New Approach to Time-Optimal Path Parameterization Based on Reachability Analysis. IEEE Transactions on Robotics 34, 645-659. \href{https://doi.org/10.1109/tro.2018.2819195}{doi:10.1109/tro.2018.2819195}.
\bibitem{prm} Kavraki, L.E., Svestka, P., Latombe, J.C., Overmars, M.H., 1996. Probabilistic roadmaps for path planning in high-dimensional configuration spaces. IEEE Transactions on Robotics and Automation 12, 566-580. \href{https://doi.org/10.1109/70.508439}{doi:10.1109/70.508439}.
\bibitem{connect} Kuffner, J.J., LaValle, S.M., 2000. RRT-connect: An efficient approach to single-query path planning. Proceedings of the IEEE International Conference on Robotics and Automation, 995-1001 vol.2. \href{https://doi.org/10.1109/robot.2000.844730}{doi:10.1109/robot.2000.844730}.
\bibitem{kino2} LaValle, S.M., Kuffner, J.J., 2001. Randomized Kinodynamic Planning. The International Journal of Robotics Research 20, 378-400. \href{https://doi.org/10.1177/02783640122067453}{doi:10.1177/02783640122067453}.
\bibitem{kpiece} \c{S}ucan, I.A., Kavraki, L.E., 2009. Kinodynamic Motion Planning by Interior-Exterior Cell Exploration. Springer Tracts in Advanced Robotics, 449-464. \href{https://doi.org/10.1007/978-3-642-00312-7_28}{doi:10.1007/978-3-642-00312-7\_28}.
\bibitem{ompl} Sucan, I.A., Moll, M., Kavraki, L.E., 2012. The Open Motion Planning Library. IEEE Robotics \& Automation Magazine 19, 72-82. \href{https://doi.org/10.1109/mra.2012.2205651}{doi:10.1109/mra.2012.2205651}.
\bibitem{rrt} Karaman, S., Frazzoli, E., 2011. Sampling-based algorithms for optimal motion planning. The International Journal of Robotics Research 30, 846-894. \href{https://doi.org/10.1177/0278364911406761}{doi:10.1177/0278364911406761}.
\bibitem{fmt} Janson, L., Schmerling, E., Clark, A., Pavone, M., 2015. Fast marching tree: A fast marching sampling-based method for optimal motion planning in many dimensions. The International Journal of Robotics Research 34, 883-921. \href{https://doi.org/10.1177/0278364915577958}{doi:10.1177/0278364915577958}.
\bibitem{informed} Gammell, J.D., Srinivasa, S.S., Barfoot, T.D., 2014. Informed RRT*: Optimal sampling-based path planning focused via direct sampling of an admissible ellipsoidal heuristic. 2014 IEEE/RSJ International Conference on Intelligent Robots and Systems, 2997-3004. \href{https://doi.org/10.1109/iros.2014.6942976}{doi:10.1109/iros.2014.6942976}.
\bibitem{bit} Gammell, J.D., Barfoot, T.D., Srinivasa, S.S., 2020. Batch Informed Trees (BIT*): Informed asymptotically optimal anytime search. The International Journal of Robotics Research 39, 543-567. \href{https://doi.org/10.1177/0278364919890396}{doi:10.1177/0278364919890396}.
\bibitem{sst} Li, Y., Littlefield, Z., Bekris, K.E., 2016. Asymptotically optimal sampling-based kinodynamic planning. The International Journal of Robotics Research 35, 528-564. \href{https://doi.org/10.1177/0278364915614386}{doi:10.1177/0278364915614386}.
\bibitem{mpnet} Qureshi, A.H., Simeonov, A., Bency, M.J., Yip, M.C., 2019. Motion Planning Networks. 2019 International Conference on Robotics and Automation (ICRA), 2118-2124. \href{https://doi.org/10.1109/icra.2019.8793889}{doi:10.1109/icra.2019.8793889}.
\bibitem{eaai1} Kicki, P., Gawron, T., \'{C}wian, K., Ozay, M., Skrzypczy\'{n}ski, P., 2021. Learning from experience for rapid generation of local car maneuvers. Engineering Applications of Artificial Intelligence 105, 104399. \href{https://doi.org/10.1016/j.engappai.2021.104399}{doi:10.1016/j.engappai.2021.104399}.
\bibitem{eaai3} Fehér, Á., Domina, Á., Bárdos, Á., Aradi, S., Bécsi, T., 2025. Path planning via reinforcement learning with closed-loop motion control and field tests. Engineering Applications of Artificial Intelligence 142, 109870. \href{https://doi.org/10.1016/j.engappai.2024.109870}{doi:10.1016/j.engappai.2024.109870}.
\bibitem{diffusion} Carvalho, J., Le, A.T., Kicki, P., Koert, D., Peters, J., 2025. Motion Planning Diffusion: Learning and Adapting Robot Motion Planning With Diffusion Models. IEEE Transactions on Robotics 41, 4881-4901. \href{https://doi.org/10.1109/tro.2025.3593109}{doi:10.1109/tro.2025.3593109}.
\bibitem{stomp} Kalakrishnan, M., Chitta, S., Theodorou, E., Pastor, P., Schaal, S., 2011. STOMP: Stochastic trajectory optimization for motion planning. 2011 IEEE International Conference on Robotics and Automation, 4569-4574. \href{https://doi.org/10.1109/icra.2011.5980280}{doi:10.1109/icra.2011.5980280}.
\bibitem{pso} Kennedy, J., Eberhart, R., None. Particle swarm optimization. Proceedings of ICNN'95 - International Conference on Neural Networks 4, 1942-1948. \href{https://doi.org/10.1109/icnn.1995.488968}{doi:10.1109/icnn.1995.488968}.
\bibitem{de} Storn, R., Price, K., 1997. Differential Evolution -- A Simple and Efficient Heuristic for global Optimization over Continuous Spaces. Journal of Global Optimization 11, 341-359. \href{https://doi.org/10.1023/a:1008202821328}{doi:10.1023/a:1008202821328}.
\bibitem{cem} de Boer, P.T., Kroese, D.P., Mannor, S., Rubinstein, R.Y., 2005. A Tutorial on the Cross-Entropy Method. Annals of Operations Research 134, 19-67. \href{https://doi.org/10.1007/s10479-005-5724-z}{doi:10.1007/s10479-005-5724-z}.
\bibitem{eaai4} Wu, X., 2026. Integration of multiple constraint-treating approaches for uncertain automatic parking path optimization utilizing competition and cooperation-driven differential evolution. Engineering Applications of Artificial Intelligence 175, 114645. \href{https://doi.org/10.1016/j.engappai.2026.114645}{doi:10.1016/j.engappai.2026.114645}.
\bibitem{cma} Hansen, N., Ostermeier, A., 2001. Completely Derandomized Self-Adaptation in Evolution Strategies. Evolutionary Computation 9, 159-195. \href{https://doi.org/10.1162/106365601750190398}{doi:10.1162/106365601750190398}.
\bibitem{cma2} Hansen, N., Müller, S.D., Koumoutsakos, P., 2003. Reducing the Time Complexity of the Derandomized Evolution Strategy with Covariance Matrix Adaptation (CMA-ES). Evolutionary Computation 11, 1-18. \href{https://doi.org/10.1162/106365603321828970}{doi:10.1162/106365603321828970}.
\bibitem{cma3} Hansen, N., Auger, A., 2013. Principled Design of Continuous Stochastic Search: From Theory to Practice. Natural Computing Series, 145-180. \href{https://doi.org/10.1007/978-3-642-33206-7_8}{doi:10.1007/978-3-642-33206-7\_8}.
\bibitem{quality} Pugh, J.K., Soros, L.B., Stanley, K.O., 2016. Quality Diversity: A New Frontier for Evolutionary Computation. Frontiers in Robotics and AI 3. \href{https://doi.org/10.3389/frobt.2016.00040}{doi:10.3389/frobt.2016.00040}.
\bibitem{cvt} Vassiliades, V., Chatzilygeroudis, K., Mouret, J.B., 2018. Using Centroidal Voronoi Tessellations to Scale Up the Multidimensional Archive of Phenotypic Elites Algorithm. IEEE Transactions on Evolutionary Computation 22, 623-630. \href{https://doi.org/10.1109/tevc.2017.2735550}{doi:10.1109/tevc.2017.2735550}.
\bibitem{cmame} Fontaine, M.C., Togelius, J., Nikolaidis, S., Hoover, A.K., 2020. Covariance matrix adaptation for the rapid illumination of behavior space. Proceedings of the 2020 Genetic and Evolutionary Computation Conference, 94-102. \href{https://doi.org/10.1145/3377930.3390232}{doi:10.1145/3377930.3390232}.
\bibitem{memap} Cully, A., 2021. Multi-emitter MAP-elites. Proceedings of the Genetic and Evolutionary Computation Conference, 84-92. \href{https://doi.org/10.1145/3449639.3459326}{doi:10.1145/3449639.3459326}.
\bibitem{cmamae} Fontaine, M., Nikolaidis, S., 2023. Covariance Matrix Adaptation MAP-Annealing. Proceedings of the Genetic and Evolutionary Computation Conference, 456-465. \href{https://doi.org/10.1145/3583131.3590389}{doi:10.1145/3583131.3590389}.
\bibitem{deb} Deb, K., 2000. An efficient constraint handling method for genetic algorithms. Computer Methods in Applied Mechanics and Engineering 186, 311-338. \href{https://doi.org/10.1016/s0045-7825(99)00389-8}{doi:10.1016/s0045-7825(99)00389-8}.
\bibitem{runar} Runarsson, T.P., Xin Yao, 2000. Stochastic ranking for constrained evolutionary optimization. IEEE Transactions on Evolutionary Computation 4, 284-294. \href{https://doi.org/10.1109/4235.873238}{doi:10.1109/4235.873238}.
\bibitem{mesh} Audet, C., Dennis, J.E., 2006. Mesh Adaptive Direct Search Algorithms for Constrained Optimization. SIAM Journal on Optimization 17, 188-217. \href{https://doi.org/10.1137/040603371}{doi:10.1137/040603371}.
\bibitem{comp} Lewis, R.M., Torczon, V., 2000. Pattern Search Methods for Linearly Constrained Minimization. SIAM Journal on Optimization 10, 917-941. \href{https://doi.org/10.1137/s1052623497331373}{doi:10.1137/s1052623497331373}.
\bibitem{solis} Solis, F.J., Wets, R.J.B., 1981. Minimization by Random Search Techniques. Mathematics of Operations Research 6, 19-30. \href{https://doi.org/10.1287/moor.6.1.19}{doi:10.1287/moor.6.1.19}.
\bibitem{dubins} Dubins, L.E., 1957. On Curves of Minimal Length with a Constraint on Average Curvature, and with Prescribed Initial and Terminal Positions and Tangents. American Journal of Mathematics 79, 497. \href{https://doi.org/10.2307/2372560}{doi:10.2307/2372560}.
\bibitem{snap} Mellinger, D., Kumar, V., 2011. Minimum snap trajectory generation and control for quadrotors. 2011 IEEE International Conference on Robotics and Automation, 2520-2525. \href{https://doi.org/10.1109/icra.2011.5980409}{doi:10.1109/icra.2011.5980409}.
\bibitem{uav} Richter, C., Bry, A., Roy, N., 2016. Polynomial Trajectory Planning for Aggressive Quadrotor Flight in Dense Indoor Environments. Springer Tracts in Advanced Robotics, 649-666. \href{https://doi.org/10.1007/978-3-319-28872-7_37}{doi:10.1007/978-3-319-28872-7\_37}.
\bibitem{fast} Zhou, B., Gao, F., Wang, L., Liu, C., Shen, S., 2019. Robust and Efficient Quadrotor Trajectory Generation for Fast Autonomous Flight. IEEE Robotics and Automation Letters 4, 3529-3536. \href{https://doi.org/10.1109/lra.2019.2927938}{doi:10.1109/lra.2019.2927938}.
\end{thebibliography}

\end{document}

